\section{第15章 ソフトウェア工学経済}
\begin{pointbox}{この資料の主張}
\term{ソフトウェア工学経済}{Software Engineering Economics}とは、ソフトウェアに関する技術的な選択を、組織の価値・費用・リスク・戦略に結び付けて判断する考え方である。単に「安いものを選ぶ」ための知識ではない。限られた人員、時間、資金をどの案に投入すれば最も価値を生むかを考えるための知識である。
\end{pointbox}

\begin{notebox}{読み方}
専門用語は原則として日本語で示し、初出では英語を括弧内に併記する。英語名は、元資料、標準、ツール、参考文献を確認するときの手がかりとして残す。本文はだである調で統一し、初学者が前提知識なしに読めるよう、概念の目的、見分け方、実務での使い所を中心に説明する。
\end{notebox}
